% !TEX root = ../../main.tex
% ============================================================
%  Fiches PSE1 — Bilan complémentaire · v0.2
%  Fichier : src/sections/neurologie.tex
%  Rôle    : Bilan neurologique :
%              · OMSPGG (Orientation, Motricité, Sensibilité,
%                Pupilles, Glasgow, Glycémie)
%              · FAST  — Suspicion AVC
%              · Rappel Glasgow (tableau Y/V/M)
%              · Rappel EVDA (échelle Eveil/Voix/Douleur/Aréactive)
% ============================================================

\begin{tcolorbox}[
  colback=pse-bleu!4, colframe=pse-bleu,
  title={\faBrain\enspace NEUROLOGIQUE — OMSPGG},
  coltitle=white, fonttitle=\bfseries\small,
  arc=2mm, boxrule=1pt,
  top=\ifA{1.5mm}{1mm}, bottom=\ifA{1.5mm}{1mm}, left=1.5mm, right=1.5mm,
]

\acronymfield{pse-bleu}{O}{rientation}{temps, lieu, identité, mémoire}
\acronymfieldsingle{pse-bleu}{M}{otricité}{mouvement spontanée, forcée, sensibilité}
\acronymfieldsingle{pse-bleu}{S}{ensibilité}{touché, douleur, température}
\acronymfieldsingle{pse-bleu}{P}{upilles}{isocorie, réactivité, mydriase (OO), myosis ($\cdot$$\cdot$)}
\acronymfieldsingle{pse-bleu}{G}{lasgow}{Y(4) + V(5) + M(6) = /15}
\acronymfieldsingle{pse-bleu}{G}{lycémie}{N : 0,8–1,2\,g/L | 3,3–6,6\,mmol/L}

\ifA{\vspace{1mm}}{}
\blocktitle{pse-rouge}{FAST — Suspicion AVC}
\acronymfieldsingle{pse-rouge}{F}{ace}{asymétrie, déviation}
\acronymfieldsingle{pse-rouge}{A}{rm}{faiblesse, paralysie}
\acronymfieldsingle{pse-rouge}{S}{peech}{trouble parole / compréhension}
\acronymfieldsingle{pse-rouge}{T}{ime}{heure début des signes}

\ifA{\vspace{1mm}}{}
\blocktitle{pse-bleu}{Rappel Glasgow}

{\smallscript
\renewcommand{\arraystretch}{1.15}
\ifx\PaperFormat\PaperFormatA\def\glasgownum{4mm}\else\def\glasgownum{3mm}\fi
\begin{tabular}{@{}|>{\centering\arraybackslash}p{\glasgownum}|l|>{\centering\arraybackslash}p{\glasgownum}|l|>{\centering\arraybackslash}p{\glasgownum}|l|@{}}
\hline
\rowcolor{pse-bleu!15}
\multicolumn{2}{@{}|c|}{\textbf{\textcolor{pse-bleu}{Yeux (Y)}}} &
\multicolumn{2}{c|}{\textbf{\textcolor{pse-bleu}{Verbal (V)}}} &
\multicolumn{2}{c|}{\textbf{\textcolor{pse-bleu}{Moteur (M)}}} \\
\hline
4 & Spontanée       & 5 & Orientée          & 6 & Obéit aux ordres \\
\hline
3 & À la demande    & 4 & Confuse           & 5 & Orientée \\
\hline
2 & À la douleur    & 3 & Inappropriée      & 4 & Évitement \\
\hline
1 & Aucune          & 2 & Incompréhensible  & 3 & Flexion anormale \\
\hline
\cellcolor{pse-bleu} & \cellcolor{pse-bleu} & 1 & Aucune            & 2 & Extension \\
\hline
\cellcolor{pse-bleu} & \cellcolor{pse-bleu} & \cellcolor{pse-bleu} & \cellcolor{pse-bleu} & 1 & Aucune \\
\hline
\end{tabular}
}

\ifA{\vspace{1mm}}{\vspace{0.4mm}}\noindent
{\colorbox{pse-bleu}{\color{white}\bfseries\scriptsize\,Rappel EVDA\,}}\enspace
{\small
  $\square$\,\textbf{\textcolor{pse-bleu}{E}}veil\enspace
  $\square$\,\textbf{\textcolor{pse-bleu}{V}}oix\enspace
  $\square$\,\textbf{\textcolor{pse-bleu}{D}}ouleur\enspace
  $\square$\,\textbf{\textcolor{pse-bleu}{A}}réactive%
}
\par
\end{tcolorbox}
